\documentclass[10pt,a4paper]{amsart}
\usepackage[margin=19mm]{geometry}
\usepackage{amsmath,amssymb,mathtools}
\usepackage[scr=rsfs,cal=euler]{mathalfa}
\usepackage{xfrac}
\usepackage[unicode,hidelinks,bookmarks=false]{hyperref}
\newcommand{\ie}{i.e.\ }
\newcommand{\cf}{cf.\ }
\newcommand{\resp}{respectively\ }
\newcommand{\Subsection}{\subsection}
\providecommand{\nobreakdash}{\nobreak\hbox{-}}
\setlength{\emergencystretch}{2em}
\multlinegap0pt
\usepackage{mathtools}
\usepackage{amssymb}
\newtheorem{lem}{Lemma}
\newtheorem{thm}{Theorem}
\newcommand{\PrimeFactor}{\textsc{PrimeFactor}}
\newcommand{\abs}[1]{\lvert #1\rvert}
\title{Prime Factorization in Models of $\text{PV}_1$}
\author{Ondřej Ježil}
\date{Source: arXiv:2505.14516v5 (CC BY 4.0)}
\begin{document}
\maketitle

\noindent\emph{Source and licence.} Prime Factorization in Models of $\text{PV}_1$. Version arXiv:2505.14516v5 (CC BY 4.0), distributed by arXiv under the Creative Commons Attribution 4.0 International licence, http://creativecommons.org/licenses/by/4.0/, as recorded in the arXiv licence metadata for this exact version and shown on the abstract page https://arxiv.org/abs/2505.14516v5. Full licence notice: \texttt{licenses/arxiv-2505.14516-CC-BY-4.0.txt}.\par\smallskip

\noindent\textit{Editorial scope.} Counted unit: the article's thm thrmmain (source0.tex lines 202--207). The article's complete original proof of it is delivered with it (source0.tex lines 208--222, one \texttt{proof} environment of 15 lines). No mathematics is written, repaired or re-derived: the statement and the proof are the article's own lines, in the article's own notation, and no retained line is substituted or re-flowed.  Retained uncounted as the source-local context the proof needs: lines 179--183; lines 190--196.  Nothing was omitted from any retained span, so the package carries no omission record.  No retained line contains a citation, so the wrapper carries no bibliography and no bibliography entry.  No retained line contains a figure environment, an image-inclusion command or drawing code, so no image is delivered and the package declares no asset dependency.  Editorial formatting only, in addition: an AMS \texttt{amsart} preamble that restates the article's own declarations and macros used by the retained lines, namely the environments \texttt{lem}, \texttt{thm} on separate counters, together with the standard packages those lines need.  The retained environments print in this document as Lemma 1, Theorem 1 in order of appearance, whereas the article prints the same items with the numbering of its own full document.  The complete native source of the article as distributed in the arXiv e-print for this exact version is bundled unchanged as \texttt{sources/178444/source0.tex}.
\begin{lem}\label{lemmabreak}
    Assume $s$ is a polynomial-time function, $c\geq 1$, $d=2^c$ and $p_1,\dots,p_d$ are distinct primes such that for any teacher $t$ which is $\PrimeFactor_0$-correcting on the input $x=\prod_{i=1}^d x_i$, the computation of $(s,t,c)$ on the input $x$ contains some prime factor of $x$ as one of the student's answers $y_i$.
    
    Then, there are distinct indices $l,k\in \{1,\dots,d\}$, such that $s$ with $\{p_1,\dots,p_d\}$ breaks $p_lp_k$. 
\end{lem}

\begin{lem}\label{lemmprob}
    Let $\Omega$ be a set, $d\geq 1$ and $F$ a function from subsets of $\Omega$ of size $d$ such that
    \[\forall T\subseteq \Omega, \abs{T}=d: F(T)\text{ is a non-empty set of $2$-element subsets of $T$.} \]
    Then,
    \[\Pr_{x_1,\dots,x_d \sim \Omega}[\{x_1,x_2\} \in F(\{x_1,\dots,x_d\})|x_1,\dots,x_d \text{ are distinct}]\geq \binom{d}{2}^{-1},\]
    where $x_1,\dots,x_d$ are sampled uniformly and independently.
\end{lem}

\begin{thm}\label{thrmmain}
    Assume $s$ is a polynomial-time function and $c,d\geq 1$ such that for any distinct primes $p_1,\dots,p_d$, there are some $1\leq l <k \leq d$ such that $s$ with $\{p_1,\dots,p_d\}$ breaks $p_lp_k$ (during a $c$-round computation). 
    Then there is $r>0$ and a polynomial time function $f$ such that for every finite set of primes $D$ we have
    \[\Pr_{\substack{p,q \sim D\\ p_1,\dots,p_{d-2} \sim D}}  \left[f(pq,p_1,\dots,p_{d-2}) \in \{p,q\}\right]\geq r,\]
    where $p,q,p_1,\dots,p_{d-2}$ are sampled uniformly and independently.
\end{thm}

\begin{proof}


    Consider the following algorithm $f$: On the input $(pq,p_1,\dots,p_{d-2})$, it first checks whether all numbers $p_1,\dots,p_{d-2}$ are distinct and do not divide $pq$ and then it tries to simulate for every permutation $\pi$ on $\{p,q,p_1,\dots,p_{d-2}\}$ protocols between the student $s$ and the teacher $t_{(\pi(p),\pi(q),\pi(p_1),\dots,\pi(p_{d-2}))}$ on the input $x=pq\prod_{i=1}^{d-2}p_i$ in the following way:
        
        It iterates over every permutation of $\{*_1,*_2, p_1,\dots,p_{d-2}\}$, where $*_1$~and~$*_2$ are placeholder values for $p$ and $q$ which are unknown to $f$. It then simulates the communication with the teacher $t_{(\pi(p),\pi(q),\pi(p_1),\dots,\pi(p_{d-2}))}$, where $\pi$ is the induced permutation, until division by exactly one of $p$ or $q$ is needed to proceed, in which case the simulation is aborted. Importantly, if it happens that the given permutation leads to $s$ with $\{p,q,p_1,\dots,p_{d-2}\}$ breaking $p,q$, then the knowledge of either $p$ or $q$ was not needed in the simulation of the teacher, as the knowledge of the product $pq$ suffices for every answer of the teacher. 
        
        After each answer of the student $s$, the algorithm $f$ tries to take $\gcd$ of it and $pq$ and if it finds a proper divisor it returns it.

    Let $F(p_1',\dots,p_d')$ be a function which outputs the set of all pairs of primes which are broken by the student $s$ with $\{p_1',\dots,p_d'\}$. By Lemma~\ref{lemmabreak}, this set is always non-empty and thus $F$ satisfies the conditions of Lemma~\ref{lemmprob}, which implies 
    \[\Pr_{p,q,p_1,\dots,p_{d-2}\sim D}[f(pq,p_1,\dots,p_{d-2})\in\{p,q\}|p,q,p_1,\dots,p_{d-2}\text{ distinct}]\geq \binom{d}{2}^{-1}.\]

    If $\abs{D}\leq 4 \binom{d}{2}$, then the probability of $p_1$ dividing $pq$ is at least $1/(4\binom{d}{2})$, and if $\abs{D}\geq 4\binom{d}{2}$, then the probability of $p,q,p_1,\dots,p_{d-2}$ being all distinct is at least \[\prod_{i=0}^{d-1} \left(1-\frac{i}{\abs{D}}\right ) \geq \left(1- \frac{d}{4 \binom{d}{2}}\right)^{d-1} \geq (1/2),\]
    and thus, combining this with the previous paragraph the probability that $f$ finds a factor of $pq$ given $p_1,\dots, p_{d-2}$ is at least $\frac{1}{4\binom{d}{2}}$.
\end{proof}

\end{document}
